\chapter{Value Function Approximation}

\section{Introduction and Motivation}
\subsection{Introduction}
In the previous lectures we represented value functions through a lookup table:
each state $s$ (or each pair $(s,a)$) has its own stored value.
This approach is simple but does not scale to large MDPs, because:

\begin{compactitemize}
    \item the number of states can be extremely large (for example Backgammon $10^{20}$ states; Go $10^{170}$ states);
    \item the state space can be continuous and therefore not enumerable;
    \item learning the value of each state individually requires an unrealistic amount of experience.
\end{compactitemize}

We therefore resort to a \textbf{parametric approximating function} that represents the value function with a reduced number of parameters:

\[
\hat{v}(s,\mathbf{w}) \approx v^\pi(s),
\qquad
\hat{q}(s,a,\mathbf{w}) \approx q^\pi(s,a).
\]

The goal is to generalize as much as possible from the observed states to those not yet visited.

\subsection{Large-Scale Reinforcement Learning}

Reinforcement Learning is applied to high-dimensional domains, such as:

\begin{compactitemize}
    \item complex games (Backgammon, Go);
    \item robotics and continuous control;
    \item industrial dynamical systems.
\end{compactitemize}

In such scenarios it is not possible to use a table.
A compact representation and an efficient update methodology therefore become fundamental.

In our context we assume \textbf{differentiable} approximators, so as to allow updates through gradient descent.

\subsection{Which Function Approximator?}

Unlike supervised learning, the data generated by an RL agent have some critical properties:

\begin{compactitemize}
    \item they are not i.i.d.;
    \item the distribution of visited states depends on the policy, which evolves over time;
    \item the targets themselves (such as TD targets) depend on the parameters and are not stationary.
\end{compactitemize}

These characteristics cause problems for gradient-based methods.
When the data are not i.i.d., the gradient estimate can become unstable: successive descent directions are not consistent and can produce oscillations or divergence.
This phenomenon is particularly evident with non-linear approximators, and it is one of the main reasons why techniques such as replay buffers, target networks and variance control methods are introduced.
\section{Gradient-Based Function Approximation}
\subsection{Gradient Descent (Recap)}

Given a differentiable objective $J(\mathbf{w})$, its gradient is:

\[
\nabla_{\mathbf{w}} J(\mathbf{w})
=
\begin{bmatrix}
\frac{\partial J}{\partial w_1} \\
\vdots \\
\frac{\partial J}{\partial w_n}
\end{bmatrix}.
\]

The gradient descent update is:

\[
\mathbf{w}_{t+1} = \mathbf{w}_t - \alpha \nabla_{\mathbf{w}} J(\mathbf{w}_t).
\]

In Reinforcement Learning we cannot compute the exact expectation of the gradient, so we use a sample estimator (SGD), which introduces noise but is computationally lighter.

\subsection{Value Function Approximation with SGD}

The goal is to minimize the squared error between the true value function and its approximation:

\[
J(\mathbf{w})
=
\mathbb{E}\left[
\left( v^\pi(S) - \hat{v}(S,\mathbf{w}) \right)^2
\right].
\]

The SGD update is:

\[
\Delta \mathbf{w}
=
\alpha
\left( v^\pi(S) - \hat{v}(S,\mathbf{w}) \right)
\nabla_{\mathbf{w}}\hat{v}(S,\mathbf{w}).
\]

Under ideal conditions this update converges to the minimum; in RL settings, the non-independence of the data requires particular attention to stability.

\subsection{Feature Vectors}

Each state is represented by a feature vector:

\[
\mathbf{x}(S)
=
\begin{bmatrix}
x_1(S) \\
\vdots \\
x_n(S)
\end{bmatrix}.
\]

Features provide a compact representation of the state.
Typical examples include: geometric distances, economic characteristics, structured descriptions of the configuration in a game.

\subsection{Linear Value Function Approximation}

A simple and stable approach consists of using a linear combination of features:

\[
\hat{v}(S,\mathbf{w})
=
\mathbf{x}(S)^T \mathbf{w}
=
\sum_{j=1}^{n} x_j(S) w_j.
\]

The corresponding objective to minimize is:

\[
J(\mathbf{w})
=
\mathbb{E}
\left[
\left( v^\pi(S) - \mathbf{x}(S)^T \mathbf{w} \right)^2
\right],
\]

which is convex in $\mathbf{w}$.
Consequently, SGD converges to the global minimum.
The update becomes:

\[
\Delta \mathbf{w}
=
\alpha
\left( v^\pi(S) - \hat{v}(S,\mathbf{w}) \right)
\mathbf{x}(S).
\]

\subsection{Table Lookup as Linear Approximation}

The tabular representation can be seen as a special case of linear approximation using indicator features:

\[
\mathbf{x}_{\text{table}}(S)
=
\begin{bmatrix}
\mathbf{1}(S = s_1) \\
\vdots \\
\mathbf{1}(S = s_n)
\end{bmatrix}.
\]

In this case, the vector $\mathbf{w}$ directly contains the values of the states:

\[
\hat{v}(S,\mathbf{w})
=
\mathbf{x}_{\text{table}}(S)^T \mathbf{w}.
\]

Although not scalable, this representation is the reference case and the conceptual starting point for more advanced feature engineering methods.
\section{Prediction with Function Approximation}

\subsection{Monte Carlo and TD Learning with Value Function Approximation}

In a Reinforcement Learning setting the true value function $v^\pi(s)$ is not available:
the agent receives only rewards and must therefore build a \textit{learning target} from experience.
The idea common to incremental prediction methods is to replace the true value function with an appropriate target, and to use this target within a gradient-based update:

\[
\Delta \mathbf{w}
=
\alpha \left( \text{target}_t - \hat{v}(S_t,\mathbf{w}) \right)
\nabla_{\mathbf{w}} \hat{v}(S_t,\mathbf{w}).
\]

In this section we analyse the targets used by Monte Carlo, TD(0) and TD($\lambda$).

\subsection{Monte Carlo with Value Function Approximation}

In the Monte Carlo method, the target for state $S_t$ is the \textbf{return} of the episode:

\[
G_t = \sum_{k=t}^{T-1} \gamma^{k-t} R_{k+1}.
\]

The return $G_t$ is an \textbf{unbiased} estimator of $v^\pi(s)$, although characterized by high variance.
After each episode we therefore obtain a series of training pairs:

\[
\langle S_1, G_1 \rangle, \langle S_2, G_2 \rangle, \dots, \langle S_T, G_T \rangle.
\]

For a linear approximator $\hat{v}(S,\mathbf{w}) = \mathbf{x}(S)^T \mathbf{w}$, the update becomes:

\[
\Delta \mathbf{w}
=
\alpha \left( G_t - \hat{v}(S_t,\mathbf{w}) \right) \mathbf{x}(S_t).
\]

Monte Carlo evaluation converges to a local optimum, even in the presence of non-linear approximation.

\subsection{TD(0) with Value Function Approximation}

The TD(0) method uses a target built \textit{online}, based on the next estimate of the value function.
The TD-target is:

\[
R_{t+1} + \gamma \hat{v}(S_{t+1}, \mathbf{w}).
\]

This target is \textbf{biased} with respect to $v^\pi(s)$, since it uses an approximate estimate, but it has lower variance and allows step-by-step updates.
The training data take the form:

\[
\langle S_1, R_2 + \gamma \hat{v}(S_2,\mathbf{w}) \rangle,
\dots,
\langle S_{T-1}, R_T \rangle.
\]

For the linear case, the update after each step is:

\[
\Delta \mathbf{w}
=
\alpha \left( R_{t+1} + \gamma \hat{v}(S_{t+1},\mathbf{w}) - \hat{v}(S_t,\mathbf{w}) \right)
\mathbf{x}(S_t)
=
\alpha \, \delta_t \, \mathbf{x}(S_t),
\]

where $\delta_t$ is the TD error.

For linear approximators, TD(0) converges close to the global optimum (the precision depends on the discount factor $\gamma$).

\subsection{TD($\lambda$) with Value Function Approximation}

The TD($\lambda$) method combines Monte Carlo and TD(0) through the parameter $\lambda \in [0,1]$.
It uses targets based on the \textit{$n$-step returns}, obtaining the \textbf{$\lambda$-return}:

\[
G_t^\lambda = (1-\lambda)
\sum_{n=1}^{\infty} \lambda^{n-1} G_t^{(n)}.
\]

The $\lambda$-return is also a \textbf{biased} estimator of the value function, but it generates on average more stable updates.
The training data are:

\[
\langle S_1, G_1^\lambda \rangle, \langle S_2, G_2^\lambda \rangle, \dots.
\]

\subsubsection{Forward View TD($\lambda$)}

The incremental update in the forward view is:

\[
\Delta \mathbf{w}
=
\alpha \left( G_t^\lambda - \hat{v}(S_t,\mathbf{w}) \right)
\nabla_{\mathbf{w}} \hat{v}(S_t,\mathbf{w})
=
\alpha \left( G_t^\lambda - \hat{v}(S_t,\mathbf{w}) \right)
\mathbf{x}(S_t).
\]

\subsubsection{Backward View TD($\lambda$)}

The backward view provides an equivalent but computationally more efficient formulation.
It introduces the vector of \textbf{eligibility traces}:

\[
E_t = \gamma \lambda E_{t-1} + \mathbf{x}(S_t),
\]

and the TD error:

\[
\delta_t = R_{t+1} + \gamma \hat{v}(S_{t+1},\mathbf{w}) - \hat{v}(S_t,\mathbf{w}).
\]

The parameter update is:

\[
\Delta \mathbf{w}
=
\alpha \, \delta_t \, E_t.
\]

The backward view is the one commonly used in implementations, since it requires neither waiting for the end of the episode nor storing multiple returns.
\section{Action-Value Function Approximation and Control}
\subsection{Action-Value Function Approximation}

The goal is not only to approximate the value function $v^\pi(s)$, but also the action-value function:

\[
\hat{q}(S,A,\mathbf{w}) \approx q^\pi(S,A).
\]

As in the case of the value function, an objective function based on the mean squared error is defined:

\[
J(\mathbf{w})
=
\mathbb{E}_\pi \left[
\left( q^\pi(S,A) - \hat{q}(S,A,\mathbf{w}) \right)^2
\right].
\]

The goal is to minimize this quantity through incremental gradient descent updates.
The stochastic gradient estimator provides the update:

\[
\Delta \mathbf{w}
=
\alpha \left( q^\pi(S_t,A_t) - \hat{q}(S_t,A_t,\mathbf{w}) \right)
\nabla_{\mathbf{w}} \hat{q}(S_t,A_t,\mathbf{w}).
\]

Since $q^\pi(s,a)$ is not directly observable, in RL methods this value is replaced by appropriate targets (Monte Carlo, TD, $\lambda$-return).

\subsection{Linear Action-Value Function Approximation}

To represent the action-value function through a linear approximator, joint state-action features are introduced:

\[
\mathbf{x}(S,A)
=
\begin{bmatrix}
x_1(S,A) \\
\vdots \\
x_n(S,A)
\end{bmatrix}.
\]

The estimated function takes the form:

\[
\hat{q}(S,A,\mathbf{w})
=
\mathbf{x}(S,A)^T \mathbf{w}
=
\sum_{j=1}^{n} x_j(S,A) w_j.
\]

The stochastic gradient descent update becomes:

\[
\nabla_{\mathbf{w}} \hat{q}(S_t,A_t,\mathbf{w}) = \mathbf{x}(S_t,A_t),
\]

\[
\Delta \mathbf{w}
=
\alpha \left( q^\pi(S_t,A_t) - \hat{q}(S_t,A_t,\mathbf{w}) \right)
\mathbf{x}(S_t,A_t).
\]

\subsection{Incremental Control Algorithms}

Analogously to policy evaluation, to learn the action-value function $q^\pi$ it is necessary to replace $q^\pi(S_t,A_t)$ with a target obtained from experience.
Depending on the method, Monte Carlo returns, TD targets, or $\lambda$-returns are used.

\subsubsection{Monte Carlo Control}

In the MC case, the target is the return:

\[
G_t = R_{t+1} + \gamma R_{t+2} + \cdots + \gamma^{T-t-1} R_T.
\]

The update becomes:

\[
\Delta \mathbf{w}
=
\alpha \left( G_t - \hat{q}(S_t,A_t,\mathbf{w}) \right)
\nabla_{\mathbf{w}} \hat{q}(S_t,A_t,\mathbf{w}).
\]

\subsubsection{TD(0) Control}

The TD(0) target for the action-value function is:

\[
R_{t+1} + \gamma \hat{q}(S_{t+1},A_{t+1},\mathbf{w}).
\]

The update is:

\[
\Delta \mathbf{w}
=
\alpha \left( R_{t+1} + \gamma \hat{q}(S_{t+1},A_{t+1},\mathbf{w}) -
\hat{q}(S_t,A_t,\mathbf{w}) \right)
\nabla_{\mathbf{w}} \hat{q}(S_t,A_t,\mathbf{w}).
\]

Defining the TD error for the action-value function:

\[
\delta_t
=
R_{t+1} + \gamma \hat{q}(S_{t+1},A_{t+1},\mathbf{w})
-
\hat{q}(S_t,A_t,\mathbf{w}),
\]

the update is written as:

\[
\Delta \mathbf{w} = \alpha \, \delta_t \, \nabla_{\mathbf{w}} \hat{q}(S_t,A_t,\mathbf{w}).
\]

\subsubsection{TD($\lambda$) Control}

For TD($\lambda$), the target is the action-value $\lambda$-return:

\[
q_t^\lambda.
\]

The forward view update is:

\[
\Delta \mathbf{w}
=
\alpha \left( q_t^\lambda - \hat{q}(S_t,A_t,\mathbf{w}) \right)
\nabla_{\mathbf{w}} \hat{q}(S_t,A_t,\mathbf{w}).
\]

\subsubsection{Backward View TD($\lambda$)}

The backward view form introduces the eligibility traces:

\[
E_t
=
\gamma \lambda E_{t-1}
+
\nabla_{\mathbf{w}} \hat{q}(S_t,A_t,\mathbf{w}),
\]

and uses the TD error:

\[
\delta_t
=
R_{t+1}
+
\gamma \hat{q}(S_{t+1},A_{t+1},\mathbf{w})
-
\hat{q}(S_t,A_t,\mathbf{w}).
\]

The parameter update is:

\[
\Delta \mathbf{w}
=
\alpha \, \delta_t \, E_t.
\]

This method is equivalent to the forward view but much more computationally efficient.

\section{Stability, Convergence and Batch Methods}

\subsection{Convergence of Prediction Algorithms}

In the case of policy evaluation with an approximate value function, the convergence of the various algorithms depends on:

\begin{compactitemize}
    \item the type of method (MC, TD(0), TD($\lambda$), etc.);
    \item the type of approximator (linear or non-linear);
    \item the on-policy vs off-policy scenario.
\end{compactitemize}

In general:

\begin{compactitemize}
    \item \textbf{Monte Carlo} with an approximate function (even non-linear) corresponds to a supervised regression problem with fixed targets (the returns $G_t$). In this case there are good guarantees of convergence to a local optimum of the error function;
    \item \textbf{TD(0)} with a \textbf{linear} approximator and \textbf{on-policy} converges to the projected fixed point of the Bellman operator (convergence ``close'' to the optimum, depending on $\gamma$);
    \item \textbf{TD(0)} (and TD($\lambda$)) with a \textbf{non-linear} approximate function or in an \textbf{off-policy} setting have in general no convergence guarantees: they may converge in practice, but there are counterexamples in which they diverge.
\end{compactitemize}

In the slides, the notation ``$x$'' indicates precisely the cases in which \textbf{there are no theoretical convergence guarantees}, even though the algorithm may behave well empirically.

\subsection{Gradient TD Learning}

Standard TD algorithms (such as TD(0), TD($\lambda$)) do not correspond to the gradient descent of an explicit objective function.
In particular, the TD update:

\[
\Delta \mathbf{w}
=
\alpha
\left( R_{t+1} + \gamma \hat{v}(S_{t+1},\mathbf{w}) - \hat{v}(S_t,\mathbf{w}) \right)
\nabla_{\mathbf{w}} \hat{v}(S_t,\mathbf{w})
\]

is not in general the gradient of a functional $J(\mathbf{w})$.

This is the reason why:

\begin{compactitemize}
    \item TD can diverge with a non-linear approximate function;
    \item TD can diverge in an off-policy scenario, even with a linear approximator.
\end{compactitemize}

\textbf{Gradient TD} algorithms were introduced precisely to solve this problem:
they approximate the true gradient descent of a well-defined error measure (for example the \emph{projected Bellman error}), introducing auxiliary weights and specifically designed updates.
In this way they recover convergence guarantees in settings where standard TD can diverge.

\subsection{Convergence of Control Algorithms}

For control (learning an optimal policy), the situation is even more delicate:

\begin{compactitemize}
    \item with a tabular representation and on-policy algorithms (e.g. tabular SARSA) there are guarantees of convergence to the optimal policy (under suitable conditions on exploration and step-size);
    \item with a \textbf{linear} approximate function and on-policy algorithms some guarantees can be obtained, but much weaker ones;
    \item with a \textbf{non-linear} approximate function (e.g. deep networks) and in particular in off-policy settings (e.g. Q-learning + replay) there are no general theoretical convergence guarantees: the behaviour is driven mainly by empirical experience and by stabilization techniques.
\end{compactitemize}

For this reason, the design of stable control algorithms with function approximation is an active research area.

\subsection{Batch Reinforcement Learning}

Incremental gradient descent algorithms update $\mathbf{w}$ step by step during the interaction.
These methods are simple to implement, but \textbf{not very sample-efficient}: each transition is used only once and then discarded.

\textbf{Batch} methods take a different perspective:

\begin{compactitemize}
    \item a fixed dataset of experience is collected;
    \item we look for the value function that best approximates the targets on that dataset.
\end{compactitemize}

In other words, given a set of \texttt{<state, target>} pairs obtained from experience, we look for the function $v(s,\mathbf{w})$ that \textbf{best fits} these data in the least-squares sense.

\subsection{Least Squares Prediction}

Suppose we have an approximator $ \hat{v}(s,\mathbf{w}) \approx v^\pi(s)$ and a dataset of experience:

\[
\mathcal{D} = \{ \langle s_1, v_1^\pi \rangle, \langle s_2, v_2^\pi \rangle, \dots, \langle s_T, v_T^\pi \rangle \},
\]

where $v_t^\pi$ is a target for the value in $s_t$ (ideally $v^\pi(s_t)$, or an estimate such as the Monte Carlo return or a $\lambda$-return).

The question is: \textit{which parameters $\mathbf{w}$ produce the best approximation of these targets?}

Least-squares algorithms look for the vector $\mathbf{w}$ that minimizes:

\[
L_{LS}(\mathbf{w})
=
\sum_{t=1}^{T} \left( v_t^\pi - \hat{v}(s_t,\mathbf{w}) \right)^2
=
\mathbb{E}_{\mathcal{D}} \left[ \left( v^\pi - \hat{v}(s,\mathbf{w}) \right)^2 \right],
\]

where the expectation is computed with respect to the empirical distribution of the data in the dataset $\mathcal{D}$.

\section{Stochastic Gradient Descent with Experience Replay}

A simple way to approximate the least-squares solution, reusing past experience, is the use of \textbf{experience replay}:

\begin{compactitemize}
    \item a dataset $\mathcal{D} = \{ \langle s_t, v_t^\pi \rangle \}$ is built from the observed transitions (for example using MC returns, TD($\lambda$), etc. as targets);
    \item \texttt{<state, target>} pairs are randomly sampled from the dataset to break the temporal correlations;
    \item an SGD step is applied to each sample.
\end{compactitemize}

Formally:

\begin{compactitemize}
    \item we draw $\langle s, v^\pi \rangle \sim \mathcal{D}$;
    \item we update
    \[
    \Delta \mathbf{w}
    =
    \alpha \left( v^\pi - \hat{v}(s,\mathbf{w}) \right)
    \nabla_{\mathbf{w}} \hat{v}(s,\mathbf{w}).
    \]
\end{compactitemize}

By repeating this procedure, SGD converges to the least-squares solution:

\[
\mathbf{w}^\pi = \arg\min_{\mathbf{w}} L_{LS}(\mathbf{w}).
\]

The central idea is that experience replay:

\begin{compactitemize}
    \item reuses the same transitions many times;
    \item reduces the correlation between samples, bringing the setting closer to the i.i.d. one;
    \item stabilizes the training of the approximating function.
\end{compactitemize}

\subsection{Experience Replay in Deep Q-Networks (DQN)}

In the case of Deep Q-Networks (DQN), experience replay is combined with the use of \textbf{fixed Q-targets} to stabilize training.

The scheme is the following:

\begin{compactitemize}
    \item the agent executes an action $a_t$ according to an $\epsilon$-greedy policy with respect to $Q(s,a;\mathbf{w})$;
    \item the transition $(s_t, a_t, r_{t+1}, s_{t+1})$ is inserted into a replay memory $\mathcal{D}$;
    \item periodically, a mini-batch of transitions $(s,a,r,s')$ is sampled from $\mathcal{D}$;
    \item for each sample the Q-learning target is computed using a \emph{fixed} weight vector $\mathbf{w}_{old}$:
    \[
    y = r + \gamma \max_{a'} Q(s',a'; \mathbf{w}_{old});
    \]
    \item the MSE between the current Q-network $Q(s,a;\mathbf{w})$ and the target $y$ is minimized:
    \[
    \mathcal{L}(\mathbf{w})
    =
    \mathbb{E}_{(s,a,r,s') \sim \mathcal{D}}
    \left[
    \left( y - Q(s,a; \mathbf{w}) \right)^2
    \right];
    \]
    \item the parameters $\mathbf{w}$ are updated with a variant of stochastic gradient descent.
\end{compactitemize}

The use of $\mathbf{w}_{old}$ (``frozen'' parameters) for computing the target is crucial:

\begin{compactitemize}
    \item it prevents the target from changing too quickly while the weights are updated;
    \item it makes the optimization more similar to a supervised regression problem with almost fixed targets;
    \item it contributes decisively to the stability of training, especially in the presence of deep approximating functions and non-i.i.d. data.
\end{compactitemize}

By combining:

\begin{compactitemize}
    \item replay buffer (to decorrelate the data and reuse the transitions);
    \item fixed target network (to make the target more stable);
\end{compactitemize}

DQN manages to effectively apply the ideas of Q-learning in the presence of non-linear approximation through deep neural networks, mitigating many of the instability problems discussed previously.
